\section{Evaluation Order and Determinism}
\subsection{Determinism}
The semantics we looked at for SIMPL are \textit{deterministic}, meaning each program has only one possible behavior. For small-step operational semantics, this means that for every expression that can be reduced by the semantics, there is exactly one way to reduce it. 

We also made choices like evaluation order to obtain deterministic semantics. For example, for the following left-to-right evaluation order: 
\[
  \text{(OP-1) } \dfrac{(e_1, s) \to (e_1^{\prime}, s^{\prime})}{(e_1 \odot e_2, s) \to (e_1^{\prime} \odot e_2, s^{\prime})} \hspace{1cm} 
  \text{(OP-2) } \dfrac{(e_2, s) \to (e_2^{\prime}, s^{\prime})}{(v \odot e_2, s) \to (v \odot e_2^{\prime}, s^{\prime})}
\]
We can instead use a right-to-left evaluation order:
\[
  \text{(OP-1) } \dfrac{(e_2, s) \to (e_2^{\prime}, s^{\prime})}{(e_1 \odot e_2, s) \to (e_1 \odot e_2^{\prime}, s^{\prime})} \hspace{1cm} 
  \text{(OP-2) } \dfrac{(e_1, s) \to (e_1^{\prime}, s^{\prime})}{(e_1 \odot v, s) \to (e_1^{\prime} \odot v, s^{\prime})}
\]
Yet these yield the same results for all SIMPL programs that do not interact with the store.

Such programs have no \textit{side effects}, and functions like these are called \textit{pure} functions that have no side effects and have deterministic results. They can be regarded as implementations of mathematical functions. 

\begin{remark}
  Side effects include any read/write to any kind of state. 
\end{remark}

However, most general-purpose programming languages have side effects, which include things like updating memory and writing to files, leading to different results from different evaluation orders. In the context of SIMPL, the only side effects are store updates. 

Most languages are implemented with a left-to-right evaluation order. Many also implement something called \textit{short-circuiting} for boolean operators, where if the boolean value of the expression can be determined from evaluating only the left operand, the right one will not be evaluated. 

\subsection{Nondeterminism}
We can include both left-to-right and right-to-left rules in the small-step operational semantics. Then we have a nondeterministic semantics, which may lead programs to have different observable behaviors. If we allow nondeterministic semantics, then a single program could potentially be evaluated to several different values. Then the interpreter should be a relation:
\[
  \textsf{Interpreter}_L \subseteq L \times V
\]
where for a program \(e\) and value \(v\):
\[
  (e, v) \in \textsf{Interpreter}_L \iff e \text{ can evaluate to } v \text{ according to the semantics. } 
\]
There are cases where nondeterministic semantics are desirable. For example, the generation of random numbers. Consider adding the ability to generate random integers in SIMPL via the expression \(\mathbf{randInt}\). We then have:
\[
  (\mathbf{randInt}, s) \to (z, s)
\]
Note that this does not enforce randomness, i.e. the value generated by \(\mathbf{randInt}\) could always be the integer 1 or a pseudorandom number. 

Another case where nondeterministic semantics might be accepted is for \textit{concurrency}, where programs are made up of parallel compositions of different programs that run at the same time. This can speed up overall computation, yet leads to nondeterministic semantics. For example, we have:
\[
  (x \coloneqq 1 \parallel y \coloneqq 2); !y
\]
It is common to model the semantics of concurrent programs using \textit{interleaving semantics}, in which the semantics of a parallel composition of programs is defined to be the interleaving of the semantic steps of the composed programs. For example:
\[
  \text{(LEFT-STEP) } \dfrac{(e_1, s) \to (e_1^{\prime}, s^{\prime})}{(e_1 \parallel e_2, s) \to (e_1^{\prime} \parallel e_2, s^{\prime})} \hspace{1cm} 
  \text{(RIGHT-STEP) } \dfrac{(e_2, s) \to (e_2^{\prime}, s^{\prime})}{(e_1 \parallel e_2, s) \to (e_1 \parallel e_2^{\prime}, s^{\prime})}
\]
\[
  \text{(LEFT-SKIP)} (\mathbf{skip} \parallel e, s) \to (e, s) \hspace{1cm} 
  \text{(RIGHT-SKIP)} (e \parallel \mathbf{skip}, s) \to (e, s)
\]
Then the four evaluations for \((x \coloneqq 1 \parallel y \coloneqq 2); !y\), starting from the empty store \(\{\}\), allowed by the semantics are as follows:

Evaluation 1:
\[
  \begin{aligned}
    ((x \coloneqq 1 \parallel y \coloneqq 2); !y, \{\}) &\to ((\mathbf{skip} \parallel y \coloneqq 2); !y, &\{x \mapsto 1\}) \\
    &\to (y \coloneqq 2; !y, &\{x \mapsto 1\}) \\
    &\to (\mathbf{skip}; !y, &\{x \mapsto 1, y \mapsto 2\}) \\
    &\to (!y, &\{x \mapsto 1, y \mapsto 2\}) \\
    &\to (2, &\{x \mapsto 1, y \mapsto 2\}) \\
  \end{aligned}
\]

Evaluation 2:
\[
  \begin{aligned}
    ((x \coloneqq 1 \parallel y \coloneqq 2); !y, \{\}) &\to ((\mathbf{skip} \parallel y \coloneqq 2); !y, &\{x \mapsto 1\}) \\
    &\to ((\mathbf{skip} \parallel \mathbf{skip}); !y, &\{x \mapsto 1, y \mapsto 2\}) \\
    &\to (\mathbf{skip}; !y, &\{x \mapsto 1, y \mapsto 2\}) \\
    &\to (!y, &\{x \mapsto 1, y \mapsto 2\}) \\
    &\to (2, &\{x \mapsto 1, y \mapsto 2\}) \\
  \end{aligned}
\]

Evaluation 3:
\[
  \begin{aligned}
    ((x \coloneqq 1 \parallel y \coloneqq 2); !y, \{\}) &\to ((x \coloneqq 1 \parallel \mathbf{skip}); !y, &\{y \mapsto 2\}) \\
    &\to (x \coloneqq 1; !y, &\{y \mapsto 2\}) \\
    &\to (\mathbf{skip}; !y, &\{x \mapsto 1, y \mapsto 2\}) \\
    &\to (!y, &\{x \mapsto 1, y \mapsto 2\}) \\
    &\to (2, &\{x \mapsto 1, y \mapsto 2\}) \\
  \end{aligned}
\]

Evaluation 4:
\[
  \begin{aligned}
    ((x \coloneqq 1 \parallel y \coloneqq 2); !y, \{\}) &\to ((x \coloneqq 1 \parallel \mathbf{skip}); !y, &\{y \mapsto 2\}) \\
    &\to ((\mathbf{skip} \parallel \mathbf{skip}); !y, &\{x \mapsto 1, y \mapsto 2\}) \\
    &\to (\mathbf{skip}; !y, &\{x \mapsto 1, y \mapsto 2\}) \\
    &\to (!y, &\{x \mapsto 1, y \mapsto 2\}) \\
    &\to (2, &\{x \mapsto 1, y \mapsto 2\}) \\
  \end{aligned}
\]

Note that while the semantics is nondeterministic, the result of the program is deterministic regardless of how the program is evaluated. However, this is not always the case for all programs.
